\documentclass[aspectratio=169,sn-mathphys-num,spanish]{beamer}

\usepackage{fontspec}
\usepackage{unicode-math}
\usepackage{graphicx}
\usepackage[dvipsnames]{xcolor}
\usepackage[sorting=none]{biblatex}
\usepackage[spanish]{babel}
\usepackage{ragged2e}
\usepackage{mathtools}
\usepackage{array}
\usepackage{booktabs}
\usepackage{tabularx}
\usepackage{tikz}
\usetikzlibrary{shapes.geometric, arrows.meta, positioning, fit}
\usepackage{siunitx}
\sisetup{
	group-separator={.},
	group-minimum-digits=4
}
\addbibresource{../references.bib}

\usetheme[numbering=counter, progressbar=frametitle, background=light]{metropolis}

\usecolortheme{dove}
\definecolor{primary}{RGB}{88, 24, 69}
\definecolor{secondary}{RGB}{0, 168, 150}

\setbeamercovered{transparent}
\setbeamercolor{frametitle}{fg=white, bg=primary}
\setbeamercolor{progress bar}{fg=secondary}
\setbeamerfont{frametitle}{size=\large, series=\bfseries}
\setbeamerfont{title}{size=\LARGE, series=\bfseries}
\setbeamercolor{alerted text}{fg=secondary}
\setbeamercolor{block title}{fg=primary}
\setbeamercolor{block body}{bg=white}

\title{Sistemas de Manufactura Flexible (FMS)}
\subtitle{Transición Energética del Relleno Sanitario Doña Juana}
\author{Julian Avila \and Bryan Martinez \and Angie Sierra}
\institute{
	Universidad Distrital Francisco José de Caldas \\
	Programa Académico de Física
}
\date{\today}

\begin{document}

\begin{frame}
	\titlepage
\end{frame}

\begin{frame}{Contenido}
	\setbeamertemplate{section in toc}[circle]
	\tableofcontents
\end{frame}

\section{El Paradigma de Manufactura Flexible}

\begin{frame}{Definición en Dinámica de Fluidos}
	\begin{block}{Concepto Fundamental}
		Transición de un modelo de disposición de residuos a un sistema de
		producción continua. El biogás se trata como una variable de estado
		(energía química) enrutada dinámicamente según la demanda.
	\end{block}

	\vspace{0.5cm}
	\textbf{Objetivos del FMS:}
	\begin{itemize}
		\item Estabilización termodinámica del caudal.
		\item Flexibilidad de enrutamiento (Poligeneración).
		\item Maximización de la transferencia de exergía.
	\end{itemize}
\end{frame}

\section{Sistemas de Ingeniería: Actual vs. Propuesto}

\begin{frame}{Sistema Operativo Actual}
	\textbf{Infraestructura Física:}
	\begin{itemize}
		\item Red de extracción: 290 pozos, 76 km de tubería HDPE.
		\item Flujo másico: $13,500$ a $16,000 \text{ m}^3/\text{h}$.
	\end{itemize}

	\vspace{0.5cm}
	\textbf{Limitaciones Críticas:}
	\begin{itemize}
		\item \textbf{Control Reactivo:} Dependencia de válvulas manuales y
			muestreo intermitente (riesgo de dilución por $O_2$).
		\item \textbf{Disipación Térmica:} Generación eléctrica confinada a
			1.7 MW. El $\sim 90\%$ del yacimiento se incinera en antorchas.
	\end{itemize}
\end{frame}

\begin{frame}{Sistema FMS a Implementar}
	\textbf{Arquitectura Ciberfísica:}
	\begin{itemize}
		\item \textbf{Telemetría:} Nodos IIoT y actuación electroneumática local.
		\item \textbf{Metrología:} Sensores NDIR ($CH_4$) y paramagnéticos ($O_2$).
	\end{itemize}

	\vspace{0.5cm}
	\textbf{Vectores de Poligeneración:}
	\begin{itemize}
		\item \textbf{Vía Eléctrica:} 15 MW de carga firme mediante 9 motores
			Jenbacher (régimen de combustión pobre).
		\item \textbf{Vía Molecular:} Fraccionamiento mediante membranas UBE
			para inyección de biometano ($>96\%$) a red urbana.
	\end{itemize}
\end{frame}

\section{Fundamentación Física}

\begin{frame}{Metrología: Óptica y Magnetismo}
	\textbf{Espectrometría NDIR ($CH_4$):} \\
	Basada en la atenuación óptica en la banda de $3.3~\mu\text{m}$. Gobernado
	por la ley de Beer-Lambert:
	$$I = I_{0} \exp(-\epsilon c L)$$
	Donde $\epsilon$ es la absortividad del enlace carbono-hidrógeno.

	\vspace{0.5cm}
	\textbf{Susceptibilidad Paramagnética ($O_2$):} \\
	Cuantificación de oxígeno basada en su interacción con gradientes
	magnéticos. Garantiza inmunidad física a la inactivación por sulfuro
	de hidrógeno ($H_2S$).
\end{frame}

\begin{frame}{Control Predictivo Basado en Modelos (MPC)}
	Sustitución del control escalar por la minimización de un funcional de costo
	multivariable sobre los espacios topológicos de estados ($X$) y acciones ($U$):

	$$ J = \sum_{k} \|y(t+k|t) - r(t+k)\|_{Q}^{2} + \sum_{j} \|\Delta u(t+j|t)\|_{R}^{2} $$

	\vspace{0.5cm}
	\textbf{Implicación Física:}
	Los operadores $Q$ y $R$ penalizan la divergencia de la trayectoria ideal,
	aislando dinámicamente intrusiones anaeróbicas sin desestabilizar la red
	global.
\end{frame}

\begin{frame}{Fraccionamiento Molecular (Membranas UBE)}
	Basado en un modelo termodinámico de \textbf{solución y difusión} a través
	de capilares de poliimida.

	\vspace{0.5cm}
	\begin{itemize}
		\item \textbf{Permeabilidad Dieléctrica:} El $CO_2$ (polar) exhibe alta
			solubilidad en la matriz asimétrica, atravesando la barrera.
		\item \textbf{Repulsión Estérica:} El $CH_4$, dada su simetría tetraédrica
			perfecta y neutralidad electromagnética, es incapaz de penetrar la
			barrera y es retenido purificado.
	\end{itemize}
\end{frame}

\section{Bibliografía}
\begin{frame}[allowframebreaks]{Referencias}
	\printbibliography
	\nocite{*}
\end{frame}

\begin{frame}{}
	\centering
	{\LARGE Gracias por su atención.}
	\vspace{0.5cm}

	\Large ¿Preguntas?
\end{frame}

\end{document}
